\section*{Chapitre \ref{ch:pl:a-tick-store-on-open-formats} --- Une base de ticks sur des formats ouverts}

\begin{solution}{exo:pl:a-tick-store-on-open-formats:1}
$14\,478\,152 - 258\,535 \times 56 = 192$ octets d'en-tête. L'en-tête de version~2 de la semaine fait 256 octets (sa disposition, en JSON,
est plus longue), donc le fichier ferait $256 + 258\,535 \times 64 = 16\,546\,496$ octets.
\end{solution}

\begin{solution}{exo:pl:a-tick-store-on-open-formats:2}
Les octets déjà écrits ne changent jamais, et les enregistrements ont une largeur fixe, si bien que les premiers
$\lfloor (\text{taille} - \text{en-tête}) / \text{largeur} \rfloor$ enregistrements sont complets et définitifs. Le lecteur ignore l'enregistrement partiel
en fin de fichier (le writer est en plein milieu de son écriture) et le lit à son passage suivant.
\end{solution}

\begin{solution}{exo:pl:a-tick-store-on-open-formats:3}
$17.5 / 6.15 = 2.85$. Les partitions sont triées par symbole et par temps et stockées en colonnes, si bien que chaque colonne est encodée (dictionnaires,
deltas) puis compressée~; les enregistrements plats gardent chaque champ à pleine largeur, et les flux IPC ne sont pas compressés.
\end{solution}

\begin{solution}{exo:pl:a-tick-store-on-open-formats:4}
223 des 17\,281 transactions, soit 1.3~\%~: pour chacune, la jointure renvoie la cotation qu'a produite l'exécution, le niveau exécuté étant retiré.
Le mid a déjà bougé dans le sens de la transaction, si bien que celle-ci semble plus proche du mid qu'elle ne l'était~: la fourchette effective
mesurée est biaisée vers zéro, et un signe de la transaction déduit de cette cotation peut s'inverser et rendre la fourchette négative.
\end{solution}

\begin{solution}{exo:pl:a-tick-store-on-open-formats:5}
Environ 2.2 millions d'événements par seconde pendant 35 minutes, soit 2\,100 secondes~: environ 4.7 milliards d'événements sur un cœur. Le flux
d'options prévoit 311 milliards de messages par jour, soit plus de soixante fois davantage~: la compaction doit tourner en parallèle par partition sur
plus de soixante cœurs, ou de façon incrémentale pendant la journée (heure par heure), ou les deux.
\end{solution}

\begin{solution}{exo:pl:a-tick-store-on-open-formats:6}
Si le numéro de version ne change pas, un ancien lecteur lit les octets du nouveau champ comme les champs qui le suivent~: des valeurs
fausses sans signal d'erreur, la pire des défaillances. Si la version change et que la largeur augmente, un ancien lecteur refuse le fichier. \texttt{check\_evolution} signale
le nouveau champ comme inséré avant les champs existants.
\end{solution}

\begin{solution}{exo:pl:a-tick-store-on-open-formats:7}
Instrument~1~: 27\,157 enregistrements~; instrument~2~: 102\,144 enregistrements, 367\,600 actions exécutées~; instrument~0 (messages système)~: 4
enregistrements. Les lecteurs C++ et Rust donnent les mêmes chiffres sur les jeux de test, qui sont écrits par le même code Python.
\end{solution}

\begin{solution}{exo:pl:a-tick-store-on-open-formats:8}
Un plantage pendant la compaction laisse la journée ni sous sa forme intrajournalière ni complète sous sa forme historique, et une requête
lancée à ce moment-là lit un mélange des deux. Il faut compacter vers de nouveaux fichiers, les vérifier (nombre de lignes, plages de numéros de séquence), puis faire basculer les lecteurs vers eux de façon atomique (un
renommage du répertoire de partition ou un manifeste), et ne supprimer les fichiers intrajournaliers qu'ensuite.
\end{solution}

\begin{solution}{pb:pl:a-tick-store-on-open-formats:1}
\begin{enumerate}
\item 564\,162 événements, 96\,206 cotations et 17\,281 transactions.
\item Le nombre magique, la version, la largeur d'enregistrement et la disposition en JSON~; complété à 64 octets pour que chaque enregistrement commence
au même alignement dans chaque fichier.
\item Les événements plats sont la disposition du chemin de négociation, pour le rejeu et les outils C++ ou Rust~; les tables IPC portent un schéma et
des chaînes, pour les dataframes et le SQL.
\item 56 et 64 octets.
\item La session doit écrire au rythme d'arrivée des données~; le tri demande la journée entière, et la compression coûte du temps
que la capture n'a pas.
\item Partitionné par date et par symbole, trié par temps à l'intérieur~: les requêtes nomment une plage de dates et un symbole, et filtrent une fenêtre temporelle.
\item 17.5~Mo en intrajournalier, 6.15~Mo après compaction.
\item 0.12~s pour 258\,535 événements, soit environ 2.2 millions d'événements par seconde sur un cœur.
\item Environ 4.7 milliards.
\item Compacter les partitions en parallèle, compacter de façon incrémentale pendant la session, et garder la forme intrajournalière peu coûteuse à convertir.
\item La couche de requête définit une vue par table, comme l'union, par nom de colonne, des partitions historiques et de la table du
jour.
\item 12\,291, 996, 1\,388, 1\,197 et 4\,473.
\item Une exécution et le changement de cotation qu'elle provoque partagent un horodatage~; le numéro de séquence les ordonne.
\item 223 transactions~; DuckDB, Polars et la référence s'accordent sur chaque ligne grâce au numéro de séquence.
\item 16~ms, 45~ms et 0.17~s.
\item \textbf{Résultat nommé.} Les 258\,535 événements du lundi se compactent en 0.12~s (environ 2.2 millions d'événements par seconde sur un cœur), de
17.5~Mo à 6.15~Mo~; une requête sur une fenêtre de cinq jours prend 18~ms, et la jointure as-of de la semaine 16~ms dans DuckDB~; à ce rythme, la fenêtre de 35 minutes
absorbe environ 4.7 milliards d'événements, soit environ un soixante-sixième de la journée prévue du flux d'options.
\item Les anciennes journées n'ont pas de colonne de plateforme et la lisent comme nulle~; les fichiers plats portent leur version et sont mis à niveau à la lecture, avec le
champ à sa valeur par défaut.
\item Un parcours quatre fois plus rapide (1.2 contre 4.8~ms) sans décodage, pour une taille 2.8 fois plus grande.
\item La séparation intrajournalière/historique et l'écriture de fin de journée vers l'historique viennent des \omterm{def:pl:a-tick-store-on-open-formats:store}{bases de ticks}~; les formats, le moteur de requête
et les dataframes sont ouverts.
\item Une seule couche de requête qui fait correspondre les colonnes par nom sur les tables du jour et les partitions historiques.
\end{enumerate}
\end{solution}

\begin{solution}{iq:pl:a-tick-store-on-open-formats:1}
Ils veulent des dispositions opposées~: la journée en cours doit être écrite au fur et à mesure de son arrivée (ajoutée, dans l'ordre temporel, sans tri), tandis que l'historique doit
être lu par plages de jours et de symboles (partitionné, trié, compressé). La compaction convertit une forme dans l'autre une seule fois, à la
clôture, et une couche de requête masque la séparation.
\iqlookfor{Optimisation de l'écriture contre optimisation de la lecture, et la compaction entre les deux.}
\end{solution}

\begin{solution}{iq:pl:a-tick-store-on-open-formats:2}
La jointure as-of sur l'horodatage peut renvoyer la cotation que la transaction a produite elle-même, puisque les deux portent le même horodatage de bourse~: le mid
a déjà bougé et la fourchette paraît plus petite, voire négative. Il faut joindre strictement sur le numéro de séquence du flux (la dernière cotation d'un
numéro de séquence antérieur), ou sur une clé d'ordre de réception.
\iqlookfor{Égalités d'horodatage~; jointure stricte sur un ordre total.}
\end{solution}

\begin{solution}{iq:pl:a-tick-store-on-open-formats:3}
Lectures sans appels système ni copies~: les pages sont chargées au premier accès et partagées entre processus, et les enregistrements à largeur fixe sont
adressés directement. C'est une mauvaise idée pour des fichiers plus gros que la mémoire avec accès aléatoire (les défauts de page dominent), pour des fichiers qui peuvent être
tronqués sous le lecteur, et là où la latence doit être bornée (un défaut de page sur le chemin critique).
\iqlookfor{Défauts de page, partage, et les cas de défaillance.}
\end{solution}

\begin{solution}{iq:pl:a-tick-store-on-open-formats:4}
Ajouter le champ en fin d'enregistrement avec une valeur par défaut, relever la version dans les en-têtes et dans le registre, faire en sorte que chaque lecteur accepte toutes les
versions, et laisser les anciens fichiers tels quels~: les anciens lecteurs ignorent le champ, les nouveaux voient la valeur par défaut pour les anciennes journées.
\iqlookfor{Évolution par ajout seul avec versions~; pas de réécriture de l'historique.}
\end{solution}

\begin{solution}{iq:pl:a-tick-store-on-open-formats:5}
Partitionner par date (chaque requête porte une plage de dates) et, si les fichiers restent assez gros, par symbole ou par compartiment de hachage des symboles~; trier par
symbole puis par temps à l'intérieur~; dimensionner les \omterm{def:pl:columnar-formats:rowgroup}{groupes de lignes} selon la mesure. 8\,000 symboles par jour donneraient de minuscules fichiers par symbole pour les valeurs peu liquides
d'où l'intérêt de les regrouper par compartiments.
\iqlookfor{Partitions par date, ordre par symbole, et le problème des petits fichiers.}
\end{solution}

\begin{solution}{iq:pl:a-tick-store-on-open-formats:6}
Capturer dans des \omterm{def:pl:a-tick-store-on-open-formats:store}{fichiers en ajout seul} (append-only) et en Arrow IPC pendant la session~; compacter à la clôture en Parquet partitionné par date et
par symbole, trié par temps, écrit dans de nouveaux fichiers et basculé de façon atomique~; DuckDB ou Polars sur les deux, au moyen de vues qui se font correspondre par nom~;
des schémas versionnés. Mesurer d'abord le nombre d'octets par événement sous chaque forme, le rythme de compaction par rapport à la fenêtre entre la clôture et la recherche, et
la latence des requêtes principales (jointures as-of, fenêtres).
\iqlookfor{Le cycle complet avec basculement atomique, et les mesures qui le dimensionnent.}
\end{solution}
